\documentclass[a4paper,12pt,twoside]{memoir}

% Castellano
\usepackage[spanish,es-tabla]{babel}
\selectlanguage{spanish}
\usepackage[utf8]{inputenc}
\usepackage[T1]{fontenc}
\usepackage{lmodern} % Scalable font
\usepackage{microtype}
\usepackage{placeins}
\usepackage{xspace}

\RequirePackage{booktabs}
\RequirePackage[table]{xcolor}
\RequirePackage{xtab}
\RequirePackage{multirow}

% Links
\usepackage[colorlinks]{hyperref}
\hypersetup{
	allcolors = {red}
}

% Ecuaciones
\usepackage{amsmath}

% Rutas de fichero / paquete
\newcommand{\ruta}[1]{{\sffamily #1}}

% Párrafos
\nonzeroparskip


% Imagenes
\usepackage{graphicx}
\newcommand{\imagen}[2]{
	\begin{figure}[!h]
		\centering
		\includegraphics[width=0.9\textwidth]{#1}
		\caption{#2}\label{fig:#1}
	\end{figure}
	\FloatBarrier
}

\newcommand{\imagenflotante}[2]{
	\begin{figure}%[!h]
		\centering
		\includegraphics[width=0.9\textwidth]{#1}
		\caption{#2}\label{fig:#1}
	\end{figure}
}



% El comando \figura nos permite insertar figuras comodamente, y utilizando
% siempre el mismo formato. Los parametros son:
% 1 -> Porcentaje del ancho de página que ocupará la figura (de 0 a 1)
% 2 --> Fichero de la imagen
% 3 --> Texto a pie de imagen
% 4 --> Etiqueta (label) para referencias
% 5 --> Opciones que queramos pasarle al \includegraphics
% 6 --> Opciones de posicionamiento a pasarle a \begin{figure}
\newcommand{\figuraConPosicion}[6]{%
  \setlength{\anchoFloat}{#1\textwidth}%
  \addtolength{\anchoFloat}{-4\fboxsep}%
  \setlength{\anchoFigura}{\anchoFloat}%
  \begin{figure}[#6]
    \begin{center}%
      \Ovalbox{%
        \begin{minipage}{\anchoFloat}%
          \begin{center}%
            \includegraphics[width=\anchoFigura,#5]{#2}%
            \caption{#3}%
            \label{#4}%
          \end{center}%
        \end{minipage}
      }%
    \end{center}%
  \end{figure}%
}

%
% Comando para incluir imágenes en formato apaisado (sin marco).
\newcommand{\figuraApaisadaSinMarco}[5]{%
  \begin{figure}%
    \begin{center}%
    \includegraphics[angle=90,height=#1\textheight,#5]{#2}%
    \caption{#3}%
    \label{#4}%
    \end{center}%
  \end{figure}%
}
% Para las tablas
\newcommand{\otoprule}{\midrule [\heavyrulewidth]}
%
% Nuevo comando para tablas pequeñas (menos de una página).
\newcommand{\tablaSmall}[5]{%
 \begin{table}
  \begin{center}
   \rowcolors {2}{gray!35}{}
   \begin{tabular}{#2}
    \toprule
    #4
    \otoprule
    #5
    \bottomrule
   \end{tabular}
   \caption{#1}
   \label{tabla:#3}
  \end{center}
 \end{table}
}

%
% Nuevo comando para tablas pequeñas (menos de una página).
\newcommand{\tablaSmallSinColores}[5]{%
 \begin{table}[H]
  \begin{center}
   \begin{tabular}{#2}
    \toprule
    #4
    \otoprule
    #5
    \bottomrule
   \end{tabular}
   \caption{#1}
   \label{tabla:#3}
  \end{center}
 \end{table}
}

\newcommand{\tablaApaisadaSmall}[5]{%
\begin{landscape}
  \begin{table}
   \begin{center}
    \rowcolors {2}{gray!35}{}
    \begin{tabular}{#2}
     \toprule
     #4
     \otoprule
     #5
     \bottomrule
    \end{tabular}
    \caption{#1}
    \label{tabla:#3}
   \end{center}
  \end{table}
\end{landscape}
}

%
% Nuevo comando para tablas grandes con cabecera y filas alternas coloreadas en gris.
\newcommand{\tabla}[6]{%
  \begin{center}
    \tablefirsthead{
      \toprule
      #5
      \otoprule
    }
    \tablehead{
      \multicolumn{#3}{l}{\small\sl continúa desde la página anterior}\\
      \toprule
      #5
      \otoprule
    }
    \tabletail{
      \hline
      \multicolumn{#3}{r}{\small\sl continúa en la página siguiente}\\
    }
    \tablelasttail{
      \hline
    }
    \bottomcaption{#1}
    \rowcolors {2}{gray!35}{}
    \begin{xtabular}{#2}
      #6
      \bottomrule
    \end{xtabular}
    \label{tabla:#4}
  \end{center}
}

%
% Nuevo comando para tablas grandes con cabecera.
\newcommand{\tablaSinColores}[6]{%
  \begin{center}
    \tablefirsthead{
      \toprule
      #5
      \otoprule
    }
    \tablehead{
      \multicolumn{#3}{l}{\small\sl continúa desde la página anterior}\\
      \toprule
      #5
      \otoprule
    }
    \tabletail{
      \hline
      \multicolumn{#3}{r}{\small\sl continúa en la página siguiente}\\
    }
    \tablelasttail{
      \hline
    }
    \bottomcaption{#1}
    \begin{xtabular}{#2}
      #6
      \bottomrule
    \end{xtabular}
    \label{tabla:#4}
  \end{center}
}

%
% Nuevo comando para tablas grandes sin cabecera.
\newcommand{\tablaSinCabecera}[5]{%
  \begin{center}
    \tablefirsthead{
      \toprule
    }
    \tablehead{
      \multicolumn{#3}{l}{\small\sl continúa desde la página anterior}\\
      \hline
    }
    \tabletail{
      \hline
      \multicolumn{#3}{r}{\small\sl continúa en la página siguiente}\\
    }
    \tablelasttail{
      \hline
    }
    \bottomcaption{#1}
  \begin{xtabular}{#2}
    #5
   \bottomrule
  \end{xtabular}
  \label{tabla:#4}
  \end{center}
}



\definecolor{cgoLight}{HTML}{EEEEEE}
\definecolor{cgoExtralight}{HTML}{FFFFFF}

%
% Nuevo comando para tablas grandes sin cabecera.
\newcommand{\tablaSinCabeceraConBandas}[5]{%
  \begin{center}
    \tablefirsthead{
      \toprule
    }
    \tablehead{
      \multicolumn{#3}{l}{\small\sl continúa desde la página anterior}\\
      \hline
    }
    \tabletail{
      \hline
      \multicolumn{#3}{r}{\small\sl continúa en la página siguiente}\\
    }
    \tablelasttail{
      \hline
    }
    \bottomcaption{#1}
    \rowcolors[]{1}{cgoExtralight}{cgoLight}

  \begin{xtabular}{#2}
    #5
   \bottomrule
  \end{xtabular}
  \label{tabla:#4}
  \end{center}
}


















\graphicspath{ {./img/} }

% Capítulos
\chapterstyle{bianchi}
\newcommand{\capitulo}[2]{
	\setcounter{chapter}{#1}
	\setcounter{section}{0}
	\chapter*{#2}
	\addcontentsline{toc}{chapter}{#2}
	\markboth{#2}{#2}
}

% Apéndices
\renewcommand{\appendixname}{Apéndice}
\renewcommand*\cftappendixname{\appendixname}

\newcommand{\apendice}[1]{
	%\renewcommand{\thechapter}{A}
	\chapter{#1}
}

\renewcommand*\cftappendixname{\appendixname\ }

% Formato de portada
\makeatletter
\usepackage{xcolor}
\newcommand{\tutor}[1]{\def\@tutor{#1}}
\newcommand{\course}[1]{\def\@course{#1}}
\definecolor{cpardoBox}{HTML}{E6E6FF}
\def\maketitle{
  \null
  \thispagestyle{empty}
  % Cabecera ----------------
\noindent\includegraphics[width=\textwidth]{cabecera}\vspace{1cm}%
  \vfill
  % Título proyecto y escudo informática ----------------
  \colorbox{cpardoBox}{%
    \begin{minipage}{.8\textwidth}
      \vspace{.5cm}\Large
      \begin{center}
      \textbf{TFG del Grado en Ingeniería Informática}\vspace{.6cm}\\
      \textbf{\LARGE\@title{}}
      \end{center}
      \vspace{.2cm}
    \end{minipage}

  }%
  \hfill\begin{minipage}{.20\textwidth}
    \includegraphics[width=\textwidth]{escudoInfor}
  \end{minipage}
  \vfill
  % Datos de alumno, curso y tutores ------------------
  \begin{center}%
  {%
    \noindent\LARGE
    Presentado por \@author{}\\ 
    en Universidad de Burgos --- \@date{}\\
    Tutor: \@tutor{}\\
  }%
  \end{center}%
  \null
  \cleardoublepage
  }
\makeatother

\newcommand{\nombre}{Guillermo Paredes Muga} %%% cambio de comando

% Datos de portada
\title{Metrominuto}
\author{\nombre}
\tutor{nombre tutor}
\date{\today}

\begin{document}

\maketitle


\newpage\null\thispagestyle{empty}\newpage


%%%%%%%%%%%%%%%%%%%%%%%%%%%%%%%%%%%%%%%%%%%%%%%%%%%%%%%%%%%%%%%%%%%%%%%%%%%%%%%%%%%%%%%%
\thispagestyle{empty}


\noindent\includegraphics[width=\textwidth]{cabecera}\vspace{1cm}

\noindent D. nombre tutor, profesor del departamento de nombre departamento, área de nombre área.

\noindent Expone:

\noindent Que el alumno D. \nombre, con DNI dni, ha realizado el Trabajo final de Grado en Ingeniería Informática titulado título de TFG. 

\noindent Y que dicho trabajo ha sido realizado por el alumno bajo la dirección del que suscribe, en virtud de lo cual se autoriza su presentación y defensa.

\begin{center} %\large
En Burgos, {\large \today}
\end{center}

\vfill\vfill\vfill

% Author and supervisor
\begin{minipage}{0.45\textwidth}
\begin{flushleft} %\large
Vº. Bº. del Tutor:\\[2cm]
D. nombre tutor
\end{flushleft}
\end{minipage}
\hfill
\begin{minipage}{0.45\textwidth}
\begin{flushleft} %\large
Vº. Bº. del co-tutor:\\[2cm]
D. nombre co-tutor
\end{flushleft}
\end{minipage}
\hfill

\vfill

% para casos con solo un tutor comentar lo anterior
% y descomentar lo siguiente
%Vº. Bº. del Tutor:\\[2cm]
%D. nombre tutor


\newpage\null\thispagestyle{empty}\newpage




\frontmatter

% Abstract en castellano
\renewcommand*\abstractname{Resumen}
\begin{abstract}
En este primer apartado se hace una \textbf{breve} presentación del tema que se aborda en el proyecto.
\end{abstract}

\renewcommand*\abstractname{Descriptores}
\begin{abstract}
Palabras separadas por comas que identifiquen el contenido del proyecto Ej: servidor web, buscador de vuelos, android \ldots
\end{abstract}

\clearpage

% Abstract en inglés
\renewcommand*\abstractname{Abstract}
\begin{abstract}
A \textbf{brief} presentation of the topic addressed in the project.
\end{abstract}

\renewcommand*\abstractname{Keywords}
\begin{abstract}
keywords separated by commas.
\end{abstract}

\clearpage

% Indices
\tableofcontents

\clearpage

\listoffigures

\clearpage

\listoftables
\clearpage

\mainmatter
\capitulo{1}{Introducción}

El tema principal del proyecto se basa en la mejora de la movilidad de los peatones en las ciudades a la hora de transitar por ellas a pie. Se centra en la idea de Metrominuto, que consiste en un mapa sinóptico que une diferentes puntos de la ciudad en función de la distancia existente entre cada uno de ellos. Trata de promover el transito a pie dentro de las ciudades
Este proyecto ofrece la posibilidad de marcar tus propias rutas incluyendo los lugares o destinos que quieras visitar y calcula la mejor forma de recorrerlos en función de la distancia para minimizar el tiempo y recorrido empleados en ello.

//Hablar sobre la movilidad urbana y sobre el fomento del turismo y la movilidad a pie ?? Quizás algún gráfico?




\capitulo{2}{Objetivos del proyecto}

Este apartado explica de forma precisa y concisa cuales son los objetivos que se persiguen con la realización del proyecto. Se puede distinguir entre los objetivos marcados por los requisitos del software a construir y los objetivos de carácter técnico que plantea a la hora de llevar a la práctica el proyecto.

\subsection{Objetivos Principales}

\begin{itemize}
	\item Desarrollar una aplicación web en la que los diferentes usuarios puedan seleccionar diversos puntos en un mapa (ciudad) con el fin de recorrerlos de la forma más óptima posible.
	\item Visualizar los puntos seleccionados: información acerca de la ubicación, marcador en el mapa.
	\item Poder añadir puntos al recorrido cuando el usuario lo desee.
	\item Calcular el recorrido mas corto para el usuario de manera que pase por todos los puntos seleccionados.
\end{itemize}

\subsection{Objetivos Técnicos}
\begin{itemize}
	\item Desarrollar la aplicación web con \textit{Flask}.
	\item Utilizar el API de Google para usar los mapas.
	\item Utilizar Git como sistema de control de versiones.
\end{itemize}
\subsection{Objetivos a nivel personal}
\begin{itemize}
	\item Realizar una aportación al turismo derivada de una necesidad personal.
	\item Profundizar en el desarrollo de un proyecto.
\end{itemize}

\capitulo{3}{Conceptos teóricos}

Aquí se explican los conceptos teóricos en los que se basa el proyecto. Aborda diferentes aspectos sobre la teoría de grafos y cómo dibujarlos.

\section{Mapa de tránsito}
Un mapa de tránsito consiste en un mapa topológico y esquemático utilizado para mostrar trayectos y estaciones en el ámbito urbano, como puede ser el metro o el autobús. Los elementos principales de este tipo de mapas son:
\begin{itemize}
	\item Líneas de diferentes colores y grosores que indican las distintas líneas del medio de transporte en cuestión.
	\item Iconos o puntos que indican las paradas o estaciones del medio en el que se vaya a viajar.
	\item Diferentes iconología para señalar características significativas.
\end{itemize}
Como aplicación para este proyecto, las estaciones o paradas del mapa corresponderán con los puntos seleccionados por el usuario sobre el mapa, y las líneas serán las distancias correspondientes entre dichos puntos.

\subsection{Harry Beck}
Harry Beck fue un ingeniero electrónico del metro de Londres que trabajaba diseñando diagramas del circuito eléctrico, y que comenzó a diseñar un nuevo mapa de las líneas y estaciones de metro de su ciudad. El objetivo de la solución estaba claro: tenía que ser sencillo de leer para el público y que este pudiese reconocer claramente las distintas estaciones, salidas y traslados.
Realizó varias versiones antes de llegar a la que conocemos hoy en día, como por ejemplo las que podemos observar en las imágenes \ref{fig:MetroLondres1928} y \ref{fig:MetroLondres1933}.
\imagen{MetroLondres1928.png}{Mapa del metro de Londres de 1928.}\label{fig:MetroLondres1928}
\imagen{MetroLondres1933.png}{Mapa del metro de Londres de 1933.}\label{fig:MetroLondres1933}

\section{Metrominuto}
El concepto de Metrominuto surgió como resultado de diversas ideas sobre movilidad en la ciudad de Pontevedra. Este concepto hace referencia a un mapa sináptico, como si de un mapa de metro se tratase, que representa las distancias y los tiempos existentes entre los diferentes puntos de una ciudad. 
//Algún gráfico sobre movilidad urbana, medios de transporte....?
\imagen{MetrominutoPontevedra}{Metrominuto de Pontevedra}

Metrominuto no solo ofrece información de cara a la gente que quiere visitar la ciudad, si no que también fomenta caminar como medio de transporte en una ciudad, donde de una manera sencilla y curiosa nos muestra cómo llegar de un sitio a otro. Caminar, como ya sabemos, es la mejor solución para evitar el gran flujo de automóviles en el área urbana, y lo que ello conlleva: una constante emisión de elementos contaminantes. \\
En los orígenes de este sistema de movilidad se encuentra el estudio, por medio de la técnica DAFO (Debilidades, Amenazas, Fortalezas, Oportunidades):

\begin{description}
	\item[Debilidades:] Como el estado cambiante del tiempo, diferente ritmo al caminar dependiendo de las personas, y la comodidad de coger el coche para moverse.
	\item[Amenazas:] Prejuicios de la población.
	\item[Fortalezas:] Cuidado del medio ambiente, mayor salud y al reducir los desplazamientos en automóvil se produce como resultado una mayor seguridad en los pocos que haya.
	\item[Oportunidades:] Mejorar la ciudad, bienestar.
\end{description}
Estos planos no solo nos incitan a caminar, si no que también incluyen información útil acerca de líneas de autobús, estaciones de ferrocarril o de metro.

\subsection{Proceso de elaboración}
Los pasos a seguir para crear un \textit{Metrominuto} de forma manual son los siguientes:

\begin{steps}
	\item Consiste en la selección, dentro de una ciudad, de los puntos que se quieren representar en el mapa. Estos puntos pueden elegirse en función de su importancia, interés turístico o de los ciudadanos.
	\item Decidir qué ruta peatonal es la mas adecuada para unirlos.
	\item Considerar cómo se va a dibujar el mapa. Puede ser más o menos preciso respecto a la realidad cartográfica.
	\item Situar un punto central que sirva como punto de origen y de orientación para todos los usuarios.
	\item Realizar por medio de herramientas de mapas, como Google Maps en nuestro caso o los mapas de Bing, el cálculo de las distancias entre los diferentes puntos.
	\item Establecer una relación entre las distancias con el tiempo medio que lleva recorrerlas. Tenemos que tener en cuenta que toda la población no camina al mismo ritmo.
	\item Una vez establecidas las diferentes rutas, hacer un estudio sobre ellas para corregir errores que puedan surgir, así como la variación en el tiempo si el terreno no es uniforme o si las condiciones de tráfico y semáforos varía.
	\item Reflejar accidentes naturales o elementos de la ciudad como parques, costa, ríos\dots A través de elementos muy sencillos y con un código de colores al que estamos acostumbrados.
	\item Reflejar aspectos de la movilidad intermodal, es decir elementos como estaciones de metro, autobús, tren, etc. 
	\item Advertir de los espacios con condiciones adversas para personas con problemas de movilidad.
	\item Simplicidad, claridad y facilidad de lectura a la hora de dibujar el mapa.
	\item No sólo mostrar conexiones con el punto central establecido como referencia, sino que también debe aparecer información sobre la interconexión entre los diferentes puntos.
\end{steps}

El objetivo de este proyecto es automatizar este proceso recogido en el documento publicado por el Concello de Pontevedra~\cite{metrominuto}, ya que actualmente los Metrominutos existentes se realizan de esta forma. Con el proceso automatizado sería el mismo usuario quien realice su propio Metrominuto con los puntos de interés personalizados que él decida. Evitando que aparezcan puntos o información que no le resulta interesante.


\section{Teoría de grafos}
Es aquella rama de las matemáticas que, junto con la ciencia de computación se encarga de estudiar las propiedades de los grafos. Primeramente debemos saber que un grafo \textit{G=(V,E)} es un conjunto de vértices o nodos unidos por enlaces llamados arcos.
Existen varios tipos de grafos, pero en ese proyecto se han usado grafos no dirigidos para realizar las distintas operaciones con los nodos, es decir, los nodos corresponden con los puntos o lugares que selecciona el usuario y los arcos hacen referencia a la distancia existente entre ellos~\cite{wiki:tgrafos}. 

\subsection{Uso de grafos}
En este proyecto el resultado final es un mapa sinóptico formado por puntos o lugares (nodos) y el <<camino>> o distancia entre ellos (arcos), y para su generación es necesario realizar una serie de operaciones. 
//explicar proceso de elaboración o explicarlo en aspectos relevantes?


\section{Scalable Vector Graphics (SVG)}
SVG o gráficos vectoriales escalables, es un término que hace referencia a un formato de gráficos vectoriales bidimensionales bien sean estáticos o dinámicos en formato XML~\cite{wiki:xml}.
Permite tres tipos de elementos:
\begin{enumerate}
	\item Elementos geométricos vectoriales, como líneas, rectas o círculos.
	\item Imágenes de mapa de bits/digitales.
	\item Texto.
\end{enumerate}

Las ventajas que presenta SVG son que a estos elementos se les pueden aplicar diferentes estilos, agrupar o transformar bien sea antes de la compilación o dinámicamente. 

\section{Referencias}

Las referencias se incluyen en el texto usando cite \cite{wiki:latex}. Para citar webs, artículos o libros \cite{koza92}.


\section{Imágenes}

Se pueden incluir imágenes con los comandos standard de \LaTeX, pero esta plantilla dispone de comandos propios como por ejemplo el siguiente:

\imagen{escudoInfor}{Autómata para una expresión vacía}



\section{Listas de items}

Existen tres posibilidades:

\begin{itemize}
	\item primer item.
	\item segundo item.
\end{itemize}

\begin{enumerate}
	\item primer item.
	\item segundo item.
\end{enumerate}

\begin{description}
	\item[Primer item] más información sobre el primer item.
	\item[Segundo item] más información sobre el segundo item.
\end{description}
	
\begin{itemize}
\item 
\end{itemize}

\section{Tablas}

Igualmente se pueden usar los comandos específicos de \LaTeX o bien usar alguno de los comandos de la plantilla.

\tablaSmall{Herramientas y tecnologías utilizadas en cada parte del proyecto}{l c c c c}{herramientasportipodeuso}
{ \multicolumn{1}{l}{Herramientas} & App AngularJS & API REST & BD & Memoria \\}{ 
HTML5 & X & & &\\
CSS3 & X & & &\\
BOOTSTRAP & X & & &\\
JavaScript & X & & &\\
AngularJS & X & & &\\
Bower & X & & &\\
PHP & & X & &\\
Karma + Jasmine & X & & &\\
Slim framework & & X & &\\
Idiorm & & X & &\\
Composer & & X & &\\
JSON & X & X & &\\
PhpStorm & X & X & &\\
MySQL & & & X &\\
PhpMyAdmin & & & X &\\
Git + BitBucket & X & X & X & X\\
Mik\TeX{} & & & & X\\
\TeX{}Maker & & & & X\\
Astah & & & & X\\
Balsamiq Mockups & X & & &\\
VersionOne & X & X & X & X\\
} 

\lstset{basicstyle=\ttfamily,
	showstringspaces=false,
	commentstyle=\color{red},
	keywordstyle=\color{blue},
	frame=single
}

\capitulo{4}{Técnicas y herramientas}


\section{Técnicas}
En esta sección se explicarán las distintas herramientas utilizadas a lo largo del proyecto.

\subsection{Scrum}
Scrum es una metodología de desarrollo ágil la cual proporciona un marco de trabajo y desarrollo de productos. No es un solo proceso, si no que en esta metodología se aplican un conjunto de buenas prácticas y procesos para que el producto final sea de la mejor calidad posible.
El principal elemento del Scrum consiste en los llamados Sprints, que son ciclos de trabajo de una semana de duración. Este periodo sirve para producir un desarrollo o mejora del producto final. Estos sprints están marcados por dos reuniones:
\begin{itemize}
	\item Planificación: en ella se presentan los requisitos o avances que tiene que cumplir el proyecto, a la vez que se estiman los tiempos y se realiza la planificación.
	\item Reunión de revisión: entrega de los requisitos acordados en la reunión de planificación y el equipo analiza el sprint.
\end{itemize}
El uso de esta metodología, junto con las diversas reuniones que se realizan, permite que el producto final sea de mejor calidad ya que en todo momento se conoce el feedback del cliente y se pueden realizar distintos cambios incrementales a medida que avanza el proyecto.
Es una metodología pensada para el trabajo en equipo, por lo que en este proyecto se han mantenido las bases pero se ha adaptado la forma de trabajar, de manera que las reuniones han sido entre los tutores y el alumno y la fecha de la reunión de planificación del Sprint coincide con la fecha de revisión del sprint anterior.



\section{Herramientas}


\subsection{Estándar Python Enhancement Proposal 8 -- PEP8}
En este proyecto se ha seguido la guía de estilo \textit{PEP8}, guía única que define cómo debería estar escrito el código \textit{Python} y la forma de nombrar variables, funciones, clases o los comentarios del mismo.


\subsection{Entorno de desarrollo Integrado (IDE)}
Para el desarrollo del proyecto, se valoraron inicialmente dos editores:
\begin{itemize}	
	\item Visual Studio Code~\footnote{\url{https://code.visualstudio.com/}}.
	\item PyCharm~\footnote{\url{https://www.jetbrains.com/pycharm/}}.
\end{itemize}

PyCharm es un IDE desarrollado por la empresa JetBrains para el lenguaje Python. Existen dos ediciones, \textit{Community} y \textit{Professional}. Gracias al programa For Students~\footnote{\url{https://www.jetbrains.com/student/}} de JetBrains se está usando la versión \textit{Professional} en el proyecto, aunque la versión \textit{Community} podría usarse también para el desarrollo del proyecto..

Ofrece soporte para Flask de serie y soporte para desarrollo web mediante complementos, además de el control de estándares en el lenguaje (tanto código como comentarios), la repetición de fragmentos de código o sugerencias en el indexado del mismo.



\subsection{GitHub}
Para el control de versiones de este proyecto he utilizado GitHub, que es un repositorio en línea que emplea Git. De esta manera tenemos acceso en línea a los diferentes cambios de nuestro proyecto.
Git maneja los distintos archivos del proyecto como un conjunto de copias instantáneas.

\subsection{Lucidchart}
\textit{Lucidchart}~\footnote{\url{https://www.lucidchart.com}} es un espacio de trabajo visual diseñado para la elaboración de distintos tipos de diagramas y gráficos. En el caso de este proyecto, se ha usado para realizar los diagramas de casos de uso, clases, y secuencia, pero tiene una gran variedad de plantillas como vemos en la siguiente figura:
\imagen{lucidchart-types}{Lucidchart.}

Además, permite iniciar sesión con nuestro usuario de la Universidad y exportar nuestro trabajo a distintos formatos: png, jpg, pdf $\dots$

\subsection{Adobe illustrator}
\textit{Adobe illustrator}\footnote{\url{https://www.adobe.com/es/products/illustrator.html}} como editor de gráficos, en este caso, para la creación del logo de la aplicación web.
\imagen{logo-negativo}{Logo Metrominuto.}

\subsection{TeXstudio}
TeXstudio es un editor de código abierto para crear documentos \LaTeX. Posee numerosas características para su desarrollo. Además, para su correcto funcionamiento es necesario instalar MiKTeK.

\subsection{Google API -- Google Console}
En este proyecto, para la selección de los distintos puntos a recorrer por parte del usuario he empleado los mapas de Google. Google proporciona una plataforma para los desarrolladores en la que se puede encontrar una gran cantidad de documentación\footnote{\url{https://cloud.google.com/maps-platform/}}.
Para poder integrar en la aplicación web tanto los mapas como las diferentes funcionalidades que ofrecen debemos adquirir lo que llama API~Key\footnote{\url{https://developers-dot-devsite-v2-prod.appspot.com/maps/documentation/geocoding/get-api-key}}, la cual se trata de una clave <<privada>> para tener acceso a los servicios de su API. Para su obtención es necesario incluir tus datos bancarios, ya que durante el primer año el uso de los servicios es gratis y luego comienza a pagarse a partir de un determinado número de peticiones.
Una vez obtenida la clave, puede restringirse su uso para ciertas direcciones o dominios, de modo que puedes mantener el control de quien la usa. Además, no vale con conseguir una clave y ya esta, si no que para usar los diferentes servicios que proporciona Google hay que activar diferentes APIs.
Las APIs que se usan en este proyecto son:
\begin{itemize}
	\item \textbf{Maps JavaScript API}: Se utiliza en el cliente, de manera que se muestra el mapa al cargar la página y permite realizar diferentes acciones en él; tales como buscar, seleccionar puntos o moverte a traves de él. Algunas de estas acciones implican el uso de algunas de funcionalidades que proporcionan las APIs explicadas a continuación.
	\item \textbf{Geocoding API}: este API consta de dos elementos:
	\begin{itemize}
		\item Geocodificación: Consiste en convertir direcciones en coordenadas.
		\item Geocodificación inversa: Consiste en convertir coordenadas en una dirección legible.
	\end{itemize}
	\item \textbf{Places API}: este servicio devuelve como resultado de la petición toda la información acerca de un lugar.
	\item \textbf{Distance Matrix API}: este API proporciona tanto la distancia como el tiempo de viaje que hay entre una lista de orígenes y una de destinos. En otras palabras, como resultado devuelve la distancia y tiempo que hay entre cada origen y cada destino.
	\item \textbf{Directions API}: como respuesta nos devuelve las indicaciones a seguir para llegar desde el punto de inicio hasta el punto de destino. Además, puede configurarse para diferentes modos de trasporte, diferentes momentos de salida o llegada.
\end{itemize}

\subsection{Firebase -- Firebase Console}
Firebase~\footnote{\url{https://firebase.google.com/}} es una plataforma móvil creada por Google para el desarrollo de aplicaciones, disponible para IOS, Android y web. Su función principal es facilitar la creación de aplicaciones de elevada calidad de una más rápida, permitiendo añadir una base de usuarios que mejore la rentabilidad económica, además de la posibilidad de añadir funciones de estadísticas, bases de datos o mensajería.


Además, cuenta con varios planes para adecuar el precio con las funcionalidades que queremos incluir. En este proyecto he incluido sólo el servicio de autenticación, que forma parte del \textit{Plan Spark}~\footnote{\url{https://firebase.google.com/pricing}}, el cual es gratuito e incluye, entre otros, los servicios que vemos en la figura \ref{fig:firebase-price}.
\imagen{firebase-price}{Planes de Firebase.}



\subsubsection{FirebaseUI (User Interface) for Web -- Firebase Auth}
\textit{FirebaseUI}, desarrollado sobre \textit{Firebase Auth}~\footnote{\url{https://firebase.google.com/docs/auth}} es una librería de código abierto para la web que proporciona componentes simples y personalizables además de los \textit{software development kit} o SDKs, que permiten la eliminación de código repetitivo.


Proporciona soluciones de autenticación integrada de manera que con unas pocas líneas de código podemos iniciar sesión con direcciones de correo electrónico y contraseñas, números de teléfono o proveedores como Google, Facebook, GitHub, Twitter, Apple, Microsoft, Yahoo $\dots$. Además, también nos proporciona la posibilidad de registrarse en nuestra aplicación o casos de recuperación de contraseñas. en la figura \ref{label} podemos apreciar cuál sería el resultado de incluirlo en nuestra aplicación.
\imagen{firebaseui}{Inicio de sesión.}
Como mencionaba arriba, estas funcionalidades que conllevan muy pocas líneas, dotan a nuestra aplicación de seguridad y nos permite tener un control absoluto sobre que usuarios acceden a ella~\ref{fig:firebase-users}.
\imagen{firebase-users}{Control de usuarios.}



\subsection{Frameworks}
El término \textit{framework}~\cite{wiki:framework} significa entorno o marco de trabajo, y consiste en una estructura conceptual y tecnológica que puede servir de base para el desarrollo software. Incluye programas, bibliotecas y lenguajes interpretados.


En este proyecto se han usado para el desarrollo del mismo los frameworks \nameref{sub:flask}, \nameref{sub:jinja} y \nameref{sub:vue}.


\subsubsection{Flask}\label{sub:flask}
Flask~\footnote{\url{https://flask.palletsprojects.com/en/1.1.x/}} es un \textit{framework} de Python que nos permite crear aplicaciones cliente -- servidor de una manera mas sencilla, y que no impone ninguna limitación respecto a estructura del proyecto ni a los componentes que usar durante el desarrollo~\cite{grinberg2014flask}, aunque cave destacar que en este proyecto se ha seguido la estructura planteada por \textit{Azure} para poder realizar el despliegue.

\subsubsection{Jinja2}\label{sub:jinja}
Para que \textit{Flask} pueda hacer uso de los contenidos \textit{HTML} requiere de la utilización de este motor de plantillas, \textit{Jinja2}~\footnote{\url{https://jinja.palletsprojects.com/en/2.11.x/}}, que consiste en un lenguaje de que permite insertar datos procesados y texto predeterminado en una página \textit{HTML} mediante las etiquetas \verb|{{variable}}| ó \verb|{% expresión %}|.


\subsubsection{Vue.js}\label{sub:vue}
\textit{Vue}~\footnote{\url{https://vuejs.org/}} es un framework progresivo utilizado para construir interfaces de usuario. Está enfocado únicamente a la capa de visualización, por lo que resulta sencillo integrarlo con otras librerías o incluso en proyectos ya existentes, como es este caso.
Vue ofrece la posibilidad de instalar el cliente mediante \textit{Node.js}\footnote{\url{https://es.vuejs.org/v2/guide/installation.html}}, aunque también permite integrarlo directamente en la capa de visualización incluyendo en los templates del proyecto el \textit{cdn}:

\renewcommand{\lstlistingname}{Vue.js cdn}% Listing -> Algorithm
\renewcommand{\lstlistlistingname}{List of \lstlistingname s}
\begin{lstlisting}[language=html,caption={Versión de desarrollo.}]
<script
src="https://cdn.jsdelivr.net/npm/vue/dist/vue.js"
</script>
\end{lstlisting}
o con
\begin{lstlisting}[language=html,caption={Versión de producción.}]
<script
src="https://cdn.jsdelivr.net/npm/vue"
</script>
\end{lstlisting}
La versión de desarrollo ofrece distintas alertas o advertencias que permiten realizar diferentes trazas a la hora de localizar posibles errores, mientras que la versión de producción está optimizada en cuanto a tamaño y velocidad.
\\
Para poder utilizar \textit{Vue} en nuestros templates, además de incluir el \textit{cdn} debemos crear un componente\footnote{\url{https://vuejs.org/v2/guide/components.html}}, al que le añadiremos una variable. Esto se debe a que realmente la estructura y el nombre de los elementos de Vue es algo diferente a los que ya conocemos de \textit{HTML5}. También es importante tener en cuenta que Vue utiliza los mismos delimitadores que \textit{Flask}: \verb|{{variable}}|. Es por esto por lo que a nuestro componente o variable debemos añadir la línea \verb|delimiters:['[[',']]']| para cambiarlos.


\section{Bibliotecas}

\subsection{JQuery}
\textit{JQuery}~\cite{doc:jquery} es una biblioteca JavaScript que hace que la manipulación de documentos HTML, el manejo de eventos y las peticiones Ajax sean más sencillas de usar y manipular. Además, una de sus ventajas principales es que funciona en prácticamente todos los navegadores web.

\subsection{Networkx}
\textit{NetworkX} es la biblioteca por excelencia de Python para trabajar con grafos y redes. Permite crear, manipular y estudiar su estructura.
\\
Características:
\begin{itemize}
	\item Estructuras de datos para grafos simples, grafos dirigidos y multigrafos.
	\item Contiene la gran parte de algoritmos utilizados para el estudio y modificación de grafos.
	\item Los nodos pueden ser <<cualquier cosa>> como por ejemplo texto, imágenes o números.
	\item Los arcos pueden tener diferentes atributos o datos, como el peso, distancia\dots
	\item Es de código abierto.
\end{itemize}


\subsection{Tube Map - D3}
Esta biblioteca permite dibujar mapas muy similares a los mapas que hoy podemos ver en el metro, con sus estaciones, paradas e intersecciones~\cite{doc:tubemap}.
\\
Se intentó implementar en el proyecto durante dos Sprints, pero se llegó a la conclusión que no se podía generar la estructura necesaria en el archivo JSON para el dibujado del mapa. Esta estructura debía contener coordenadas con números enteros y estar ordenadas de tal forma que se indicasen las esquinas y cruces que debía haber entre las diferentes líneas, o en este caso, recorridos.

\subsection{SVGWRITE}
\textit{SVGWRITE}~\cite{doc:svgwritedocs}, como su propio nombre sugiere, sirve para la creación y dibujado de gráficos vectoriales escalables. Su principal inconveniente es que no permite leer o modificar un archivo ya existe, es decir, sólo permite su creación desde cero.

\subsection{Snap.svg}
\textit{Snap.svg}~\cite{doc:snapsvg} es una biblioteca JavaScript que permite crear gráficos vectoriales interactivos y fácilmente manipulables en el lado del cliente. Su uso es similar al uso de jQuery.


\subsection{\LaTeX}
Como herramienta para realizar la documentación se ha escogido \LaTeX{}~\cite{wiki:latex}, lenguaje orientado a la creación de documentos que presenten una alta calidad tipográfica.

\subsubsection{MiKTeK}
\textit{MiKTeK} es una distribución libre de \LaTeX que incluye tipografías compiladores y herramientas para generar bibliografías e índices.


\section{Despliegue}
La idea de plantear el proyecto como una aplicación web sugiere que ésta tiene que ser fácilmente accesible por lo usuarios, cosa que no es cierta si para poder usarla tienen que descargarse el código, configurar el proyecto y ejecutarlo en su propio dispositivo. Por esta razón, valoré desplegar la aplicación en Azure~\footnote{\url{https://azure.microsoft.com}}.

Elegí esta plataforma por que con el correo de la Universidad podemos tener acceso a un plan de estudiantes, llamado \textit{Azure for Students}~\footnote{\url{https://azure.microsoft.com/es-es/free/students/}}, que permite, además de acceso a diversos tutoriales sobre computación en la nube, desplegar nuestra aplicación de manera gratuita, aunque por supuesto también tiene varios planes de pago.


Link de la aplicación: \url{http://tfgmetrominuto.azurewebsites.net/}.
\capitulo{5}{Aspectos relevantes del desarrollo del proyecto}

En este apartado se van a recoger y explicar los aspectos más importantes del desarrollo del proyecto. Desde las implicaciones de las decisiones que se tomaron, hasta los numerosos y variados problemas a los que hubo que enfrentarse.

\subsection{Elección del proyecto}

El año pasado fui uno de los privilegiados de poder disfrutar de una beca Erasmus, en concreto con destino en la ciudad polaca de Gliwice~\cite{wiki:gliwice}.
Esto me llevó a conocer nuevas culturas, pero también a conocer nuevas ciudades, y en la mayoría de ocasiones el tiempo del que disponíamos para recorrerlas era muy breve. Por ello, al ver las posibilidades que ofrecía el resultado final de este TFG, me llamó la atención, ya que es una aplicación que de haberla tenido, nos habría ahorrado muchos desplazamientos, quizá inútiles al recorrer estas ciudades sin una ruta fija.

\subsection{Formación}

Para poder realizar el proyecto se necesitaban unos conocimientos no adquiridos sobre desarrollo web, tanto de la parte de servidor, en Flask, como la parte del cliente, en HTML, CSS y JavaScript. Además de para aprendizaje, los recursos se han usado también como material de consulta durante el desarrollo.

Para la parte del servidor se siguieron los siguientes libros y tutoriales:
\begin{itemize}
	\item Flask Web Development~\cite{grinberg2014flask}.
	\item The Flask Mega-Tutorial (2017)~\cite{grinberg-mega}.
	\item Deploy Python apps to Azure App Service on Linux from Visual Studio Code~\cite{deploy-flask-azure}.
	\item Flask Tutorial in Visual Studio Code~\cite{vscode-flask}.
\end{itemize}


A medida que se añadían nuevas herramientas al proyecto, su documentación
oficial también ha sido consultada en varias ocasiones, están disponibles en:
\begin{itemize}
	\item W3Schools Tutorials~\cite{w3schools}.
	\item Documentación de Flask~\cite{doc:flask}.
	\item Documentación de Bootstrap~\cite{doc:bootstrap}.
	\item Documentación de NetworkX~\cite{SciPyProceedings_11}.
	\item Tutorial FirebaseUI~\footnote{\url{https://github.com/firebase/FirebaseUi-Web\#installation}}.
\end{itemize}


\subsection{Calculo de grafos}\label{sub:grafos-aspectos-relevantes}
Para calcular el mapa sinóptico se han empleado grafos, junto con la librería~\nameref{sub:networkx} para realizar las operaciones que se explican a continuación:

\begin{steps}
	\item Calcular las distancias de todos con todos, grafo no dirigido al que llamaremos $\alpha$: añadir un arco entre cada par de nodos del grafo, asignando a cada nodo la posición que ocupa en el mapa y a cada arco el valor de la distancia que representa. Estos datos de distancia y posición los obtenemos de los APIs de Google de \textit{Places}~\footnote{\url{https://developers.google.com/maps/documentation/distance-matrix/start?hl=es}} y \textit{Distance Matrix}~\footnote{\url{https://developers.google.com/places/web-service/intro?hl=es}}.
	\item Eliminar un nodo: eliminamos un nodo del grafo junto con sus arcos asociados.
	\item Minimum spanning tree del grafo obtenido en el paso 2: evaluamos el camino mínimo para recorrer los nodos que nos quedan en el grafo.
	\item A cada arco obtenido en el grafo resultante del paso 3.
	\item Guardamos en una lista los votos de cada arco.
	\item Repetir pasos 2, 3 y 4 hasta que se haya eliminado cada nodo del grafo una única vez. De este modo, en la lista final de arcos tendremos el número de veces que aparece cada arco en el recorrido mínimo.
	\item Para cada valor ($\epsilon$) comprendido entre el mínimo y máximo de los votos anteriores:
	\begin{enumerate}
		\item Calculamos un subgrafo ($\beta$) con todos los nodos y aquellos arcos con un número de votos mayor o igual a $\epsilon$.
		\item Comprobamos que el subgrafo $\beta$ sea conexo: comprobamos que todos los nodos estén conectados. En el caso de que no lo estén, mediante NetworkX obtenemos los subconjuntos de $\beta$ y, recuperando el grafo inicial $\alpha$, unimos los subconjuntos por los arcos de menor coste. Esto quiere decir, que si por ejemplo nuestro grafo inicial $\alpha$ es (imagen~\ref{fig:explicación-grafo-conexo}):
		\imagen{explicación-grafo-conexo}{Grafo inicial.}
		y el resultado de comprobar si $\beta$ es conexo o no es (imagen~\ref{fig:explicación-grafo-inconexo}):
		\imagen{explicación-grafo-inconexo}{Grafo inconexo.}
		La unión entre ambos subconjuntos (nodo <<b>> con el resto del grafo) se realizaría por el arco que tenga una menor distancia, b--a.
		\item Guardar el subgrafo conexo obtenido.
	\end{enumerate}
	\item Generar el grafo SVG correspondiente a cada subgrafo, que se explica en la siguiente sección, \nameref{sub:dibujar-grafo-svg}.
\end{steps}



\subsection{Dibujado de grafos}\label{sub:dibujar-grafo-svg}
El principal objetivo de este proyecto se basa en obtener un grafo final de manera que éste sea fácilmente entendible por todos. Es por ello que a la hora de dibujar dicho grafo se plantean cuestiones y problemas como:
\begin{itemize}
	\item Dónde colocar los textos y a qué distancia de la línea.
	\item Cómo orientar esos textos.
	\item Evitar la superposición de los textos, tanto los referentes al punto como los referentes a la información del trayecto, con las líneas o puntos.
	\item Evitar la superposición de unos textos con otros.
	\item Tener en cuenta los colores a la hora de agrupar elementos.
	\item Mantener la posición de los puntos lo más similar posible a los marcados en el mapa.
	\item Reposicionar los puntos que vertical u horizontalmente sean casi similares para evitar líneas diagonales.
\end{itemize}

Estos problemas han resultado de gran complejidad en el desarrollo final del proyecto ya que, como mencionaba antes, es el resultado final de la aplicación web diseñada y programada para el TFG.

Un aspecto clave a tener en cuenta en el resultado final de de la aplicación es que, a pesar de que la unión entre dos puntos es una línea recta, el recorrido entre ambos puntos en la realidad no tiene por que ser recto. Dos puntos que parecen estar cerca el uno del otro pueden tener una distancia mayor que dos puntos que parecen más alejados. Para entender mejor esto fijémonos en la imagen~\ref{fig:ejemplo-distancia}. Vemos que los puntos relativos a el \textit{Arco de San Martín} y \textit{Mirador del Castillo} están más cerca el uno del otro, mientras que el punto relativo a \textit{La Estación} se encuentra más alejado de ellos.
\imagen{ejemplo-distancia}{Puntos en el mapa.}
Sin embargo, si nos fijamos en el mapa anterior podemos ver que para llegar desde el \textit{Arco de San Martín} al \textit{Mirador del Castillo} tenemos que dar un rodeo más largo de lo que parece, y puede ser el caso de que los puntos no se encuentren al mismo nivel (en términos de tiempo no es lo mismo hacer un kilómetro subiendo que un kilómetro bajando). Esto hace que, como vemos en la imagen~\ref{fig:ejemplo-distancia-tiempo}, el tiempo que tardamos entre estos dos puntos y el que tardamos desde el primero a \textit{La Estación}, que visualmente parece estar más lejos, sea el mismo.
\imagen{ejemplo-distancia-tiempo}{Tiempo entre los puntos.}

\subsection{Decoradores de Flask}
Los decoradores son funciones que envuelven y reemplazan a otra función, y al hacerlo, deben utilizar los valores de los argumentos de la función original a la nueva.

En Flask, cada vista o ruta es una función de Python, lo que significa que mediante decoradores podemos añadir funcionalidades adicionales. El más utilizado es \texttt{@route()}, que sirve para especificar la ruta y qué tipo de peticiones admite. 



En el desarrollo del proyecto se incluyó un decorador para comprobar, antes de acceder a cada vista, si el usuario había iniciado sesión o no mediante la comprobación de una variable de sesión, y en el caso de que no, redirigir a dicho usuario a la página de inicio. Se trata de algo similar a lo que hace algunos paquetes de Flask de validación de usuarios como Flask--Login~\cite{doc-flask-login}, pero que en este proyecto se ha hecho de forma manual porque se ha hecho la validación de usuarios a través de Firebase.



\subsection{Blueprints}
\textit{Blueprints}~\cite{doc:blueprints} o <<planos>> es un concepto usado en Flask para crear componentes y modular la aplicación. su funcionalidad principal es proporcionar un módulo para registrar funciones en la aplicación.




\capitulo{6}{Trabajos relacionados}

Como se menciona anteriormente en este proyecto, la inicial de Metrominuto proviene de una idea en Pontevedra que tenía como finalidad el fomento de el <<arte de caminar>>~\cite{wiki:metrominuto-pontevedra}.



\subsection{Artículos relacionados}

\subsubsection{Al menos 57 ciudades han copiado el Metrominuto pontevedrés~\cite{art:metromin-ciudades}}
Este artículo explica que no sólo ciudades españolas han replicado la idea del Concello de Pontevedra, si no que dicha idea se ha extendido a ciudades europeas como Florencia y Cagliari  en Italia o Poznan en Polonia, como podemos ver en la figura~\ref{fig:poznan}.
\imagen{poznan}{Metrominnuto de Poznan}

FIXME: tiro por artículos?


\subsection{Aplicaciones similares}
Pasominuto: https://pontevedraviva.com/web/uploads/arquivos/pasominuto.pdf

Conforme iba pasando el tiempo, muchas ciudades copiaron el modelo de Metrominuto que tenía Pontevedra, pero no solo aplicandolo al ámbito urbano, si no que también trasladándolo a otros. Es por esto que de la idea original surgen una serie de aplicaciones o herramientas similares que veremos a continuación.

\subsubsection{Pasominuto}
Esta alternativa a Metrominuto consiste en una guía de recorridos para pasear por <<los espacios más agradables>> de la ciudad, teniendo en cuanta sus distancias, pasos y tiempos. El artículo~\cite{art:pasominuto} menciona hasta un total de 29 recorridos diferentes.
\imagen{pasominuto}{Pasominuto.}

\subsubsection{Metropalyas}
Apllicación considera como <<hijo de metrominuto>> y que situa a Pontevedra en el punto equidistante de las mejores playas de las \textit{Rías Baixas}.
\imagen{metroplaya}{Metroplaya.}



FIXME : he encontrado un facebook de metrominuto https://www.facebook.com/metrominuto/  , que no se si puedo hacer mención de algún modo.
\capitulo{7}{Conclusiones y Líneas de trabajo futuras}

En este apartado se exponen las conclusiones obtenidas tras finalizar el proyecto, además de explicar posibles líneas de mejora futuras para el mismo.

\subsection{Conclusiones}
Tras ocho meses de trabajo en el desarrollo, además de trabajar fuera de la Universidad, considero que he cumplido los objetivos del proyecto, y que la aplicación resultante de todo el trabajo puede llegar a ser realmente útil, no solo para fomentar los desplazamientos a pie en la propia ciudad, si no que también a la hora de viajar y conocer nuevas ciudades.


Destacar la gran cantidad de conocimientos nuevos aprendidos, no solo en cuanto a programación, si no también en cuando a herramientas y capacidad de reunir información. He intentado usar funcionalidades y herramientas nuevas, como los frameworks y librerías explicados en el apartado \nameref{herramientas}. Creo que el resultado final es funcional, útil, y como mencionaba, me ha permitido adquirir nuevos conocimientos.


Por último, después de varios años en la Universidad, ser capaz de desarrollar un proyecto desde su inicio y que el resultado final sea una aplicación entregable y funcional es un gran logro, y personalmente, estoy muy orgulloso de ello.


\subsection{Líneas de trabajo futuras}\label{futuras}
Como futuras evoluciones del proyecto, se plantean las siguientes ideas:

\begin{itemize}
	\item Mejorar la seguridad de la aplicación: ocultar en la consola del navegador las claves mediante algún sistema de seguridad.
	\item Mejorar la tolerancia a fallos: mejorar la respuesta de la aplicación ante fallos o errores.
	\item Añadir una base de datos: con esta mejora podríamos conseguir personalizar la aplicación para cada usuario, de tal modo que este pueda guardar y editar en cualquier momento sus metrominutos.
	\item Mejorar la apariencia y maquetación de la aplicación: mejorar los estilos de la aplicación.
	\item Añadir decoración al mapa sinópico como estaciones de transporte público, parques, lagos, ríos o zonas de costa, en el caso de que existan.
	\item Añadir geolocalización (provista por otro API de Google, Geolocation\footnote{\url{https://developers-dot-devsite-v2-prod.appspot.com/maps/documentation/geolocation/intro}}): de este modo el usuario puede identificar fácilmente en qué punto del recorrido se encuentra.
	\item Permitir al usuario en el SVG agrupar zonas, dibujar parques o añadir elementos que considere interesantes.
	\item Tras la modificación de la posición de textos o puntos por parte del usuario, dar la opción al usuario de recuperar el <<magnetismo>> entre textos y nodos manteniendo la posición actual de los mismo, de tal modo que si vuelve a mover el punto, los textos se muevan junto a él.
\end{itemize}



\bibliographystyle{plain}
\bibliography{bibliografia}

\end{document}
